\subsection{Linear jet interpolation and subexponential tensor levels}

\begin{theorem}[DP56: quantitative tensor levels for sharp finite jets]\label{DP56}
Let $X$ be proper and smooth over a complete DVR $R$, let $e$ be a section with ideal $I$, and let $\mathcal L$ be relatively ample.
There are fixed integers $r\ge1$ and $b\ge0$ such that, for every $n\ge0$ and $\kappa\ge rn+b$,
\begin{equation}\label{eq:linearjetthreshold}
 H^0(X,\mathcal L^{\otimes\kappa})
 \twoheadrightarrow
 H^0\bigl(V(I^{n+1}),\mathcal L^{\otimes\kappa}|_{V(I^{n+1})}\bigr).
\end{equation}

For the actual theta line of DP54 at $p=11$, the sharp ordinary primitive-line sequences of DP46 and the first-row two-column harmonic sequence of DP51 can be realized by global integral sections of powers $\mathcal L^{\otimes\kappa_\nu}$ with
\[
 \kappa_\nu=M_\nu^2,\qquad p\nmid M_\nu,
 \qquad \log M_\nu=o(n_\nu),
 \qquad \log\kappa_\nu=o(n_\nu),
\]
where $n_\nu$ is the selected derivative degree.
The normalized selected values, including $M_\nu^{-n_\nu}$, have exactly the valuations of the corresponding pure primitive-kernel values at exponent $-1$.
The same tensor powers can be retained for number-field sections integral at the chosen prime, preserving those finite valuations.
These bounds concern tensor levels, not the archimedean sizes of the sections or their membership in a prescribed translated family.
\end{theorem}

\begin{proof}
\emph{A fixed integral projective coordinate system.}
Choose $r\ge1$ and an immersion
\[
 \iota:X\longrightarrow\mathbf P^N_R,
 \qquad \mathcal L^{\otimes r}\cong\iota^*\mathcal O(1).
\]
This is the known ample-power result, Stacks Project Lemma~29.40.4, Tag~01VT; properness makes the immersion closed.
Let $h_0,\ldots,h_N$ be the coordinate sections.
After choosing a generator of $e^*\mathcal L^{\otimes r}$, some $h_0(e)$ is a unit; rescale it to one.
Replace $h_i$ by $h_i-h_i(e)h_0$ for $i>0$.
The affine ratios $h_i/h_0$ now vanish along $e$ and generate its conormal module, because the chart is closed in affine space.
This conormal module is free of rank $g=\dim(X/R)$.
Choose $g$ ratios whose cotangent matrix has a unit minor, by choosing a basis modulo the maximal ideal and using Nakayama's lemma.
Relabel their sections $h_1,\ldots,h_g$ and put $t_i=h_i/h_0$.
Their linear matrix is invertible over $R$, so formal inversion over $R$ identifies the formal neighbourhood with $R[[t_1,\ldots,t_g]]$.

In the frame $h_0^n$, any integral $n$-jet is a polynomial
\[
 Q(t)=\sum_{|\alpha|\le n}q_\alpha t^\alpha,
 \qquad q_\alpha\in R.
\]
It is the restriction of the global section
\[
 \sum_{|\alpha|\le n}q_\alpha
 h_0^{n-|\alpha|}h_1^{\alpha_1}\cdots h_g^{\alpha_g}
 \in H^0(X,\mathcal L^{\otimes rn}).
\]
All exponents are nonnegative and every coefficient is integral.

Choose $b$ so that $\mathcal L^{\otimes m}$ is globally generated for every $m\ge b$; this is Stacks Project Proposition~28.27.13(7), Tag~01Q3.
For $\kappa\ge rn+b$, a global section $\eta$ of $\mathcal L^{\otimes(\kappa-rn)}$ has unit value at $e$.
Express the desired jet in the frame $\eta h_0^n$ and multiply the displayed homogenization by $\eta$.
This proves \eqref{eq:linearjetthreshold}, including $n=0$.\footnote{The exact known projective facts used here are \url{https://stacks.math.columbia.edu/tag/01VT}, parts~(2) and~(5), and \url{https://stacks.math.columbia.edu/tag/01Q3}, part~(7).}
In particular the inverse unit correction $U^{-\kappa}$ in DP54 is an allowed target. The change from its specified integral frame to this projective frame is an integral unit, so the resulting normalized germ agrees with the pure exponent-$-\kappa$ kernel through degree $n$.

\emph{A coefficientwise estimate in the tensor parameter.}
In the mixed-characteristic application put $K=\operatorname{Frac}R$ and $s=1/(p-1)$, with $v_p(p)=1$.
Suppose a zero-constant series $\gamma(\mathbf Z)$ in $D$ additive variables satisfies
\begin{equation}\label{eq:tensorlogprofile}
 \inf_\alpha\bigl(v_p([\mathbf Z^\alpha]\gamma)+w|\alpha|\bigr)\ge s
\end{equation}
for a rational $w>0$, and put $c=s-w$.
For a nonzero integral homogeneous polynomial $P$ of degree $n$, set
\[
 F_P(T)=P(\partial)\exp(-T\gamma)(0)\in K[T],
 \qquad
 \sigma(P)=s\max_{\alpha\in\operatorname{supp}P}
                      \sum_{i=1}^D s_p(\alpha_i).
\]
Each coefficient of this polynomial satisfies
\begin{equation}\label{eq:tensorcoefficientbound}
 v_p([T^m]F_P)\ge cn-\sigma(P),
 \qquad
 \sigma(P)\le D(1+\lfloor\log_p n\rfloor)\quad(n\ge1).
\end{equation}

To prove the coefficient assertion, choose a finite extension containing $\varpi$ with $v_p(\varpi)=w$ and set
$\widetilde\gamma(\mathbf Z)=\gamma(\varpi\mathbf Z)$.
Every coefficient of $\widetilde\gamma$ has valuation at least $s$.
For every $m\ge1$, every coefficient of $\widetilde\gamma^m/m!$ has valuation at least
\[
 ms-v_p(m!)=s\,s_p(m)\ge0.
\]
Thus every coefficient in the additive variables of
$\exp(-T\widetilde\gamma)$ is a polynomial in $T$ with integral coefficients.
Each is a polynomial because $\widetilde\gamma$ has zero constant term.
It follows that every $\mathbf Z^\alpha T^m$ coefficient of $\exp(-T\gamma)$ has valuation at least $-w|\alpha|$.
Applying $P(\partial)$ multiplies this coefficient by the integral coefficient $P_\alpha$ and by $\alpha!$.
Legendre's identity gives
\[
 v_p(\alpha!)=s|\alpha|-s\sum_i s_p(\alpha_i),
\]
which proves the first inequality in \eqref{eq:tensorcoefficientbound}.
The digit estimate proves the second.
This is a coefficientwise polynomial bound; it has not been inferred from evaluations at integral $T$.

\emph{The actual logarithm profiles.}
For a primitive logarithmic line, the complete initial-form computation of DP46 gives
\[
 \inf_n\bigl(v_p([t^n]\beta(ut,vt))+w_kn\bigr)=s.
\]
Use $\gamma(t)=\beta(ut,vt)$, $D=1$ and $P(Z)=Z^n$.
For the first-row harmonic family, set
\[
 w_*=\frac{p^{h-1}}{(p-1)(p^{h-1}+1)},
 \qquad c_*=\frac1{(p-1)(p^{h-1}+1)}.
\]
At the $j=1$ endpoint profile of DP46, the scalar and first dual weights equal $w_*$; the later dual weights are
\[
 \omega_i=s-p^{i-1}(s-w_*)<w_*\quad(i>1).
\]
The contracted inverse logarithms are integral and the group-coordinate primitive logarithm has coefficient infimum $s$.
Giving every additive input the common weight $w_*$ supplies nonnegative extra factors in the later coordinates.
Composition with the integral inverse maps therefore proves \eqref{eq:tensorlogprofile} for the full primitive logarithm.
This is the logarithmic estimate in DP51's full-polydisc corollary, before exponentiation.
For two copies, their sum also satisfies it. Take $D=2h$, that sum as $\gamma$, and the first-row polynomial $P_{a,b}$ of DP51.
Arbitrary unit arrows retain the same coefficient valuations and the unit initial coefficient in these profile computations.

\emph{Strict integer congruences and a bounded choice of level.}
Fix one selected nonzero value of degree $n\ge1$ and put
\[
 V=v_p(F_P(1)),\qquad L=cn-\sigma(P),\qquad
 B=\max\{1,\lfloor V-L\rfloor+1\}.
\]
The coefficientwise bound implies $V\ge L$, so $B>V-L$ and $B\le V-L+1$.
For any integer $T\equiv1\pmod{p^B}$,
\[
 v_p(F_P(T)-F_P(1))\ge B+L>V,
\]
because $T^m-1$ is divisible by $p^B$ for each positive integer $m$.
Hence $v_p(F_P(T))=V$.
The strict floor-plus-one choice covers integral and nonintegral $V-L$ without an endpoint error.

Let $r,b$ be the fixed interpolation constants and define
\[
 m_0=\lceil\sqrt{rn+b}\rceil,\qquad
 q_0=\left\lceil\frac{m_0-1}{p^B}\right\rceil,
 \qquad M=1+q_0p^B.
\]
Here $n\ge1$ and $r\ge1$, so $m_0\ge1$ and $q_0\ge0$.
These definitions give
\[
 M\equiv1\pmod{p^B},\qquad
 M\ge m_0,\qquad
 M<m_0+p^B\le\sqrt{rn+b}+1+p^B.
\]
Set $\kappa=M^2$. Then $\kappa\ge rn+b$, so the integral interpolation realizes the pure exponent-$-\kappa$ jet, and $\kappa\equiv1\pmod{p^B}$ preserves its selected valuation $V$.
Since $p\nmid M$, the source factor $M^{-n}$ preserves it as well.

Along a sharp sequence, $V_\nu/n_\nu\to c$ and the fixed number of variables gives $\sigma(P_\nu)=O(\log n_\nu)$.
Consequently
\[
 0\le\frac{B_\nu}{n_\nu}
 \le\frac{V_\nu-cn_\nu+\sigma(P_\nu)+1}{n_\nu}
 \longrightarrow0.
\]
The right side is nonnegative by the coefficient bound, even when $V_\nu-cn_\nu$ is negative.
It follows that
\[
 0\le\frac{\log M_\nu}{n_\nu}
 \le\frac{\log(\sqrt{rn_\nu+b}+1+p^{B_\nu})}{n_\nu}
 \longrightarrow0.
\]
Multiplying by two proves $\log\kappa_\nu=o(n_\nu)$.

\emph{Number-field sections at unchanged powers.}
For each fixed selected $\kappa_\nu$, the flat-base-change and density argument of DP54 identifies the completed global-section module with the completion of the finite free lattice over $\mathcal O_{\Q(\zeta_7),\mathfrak p}$.
Approximate the exact integral interpolant sufficiently closely to preserve its selected nonzero finite valuation.
This does not change $\kappa_\nu$, so the tensor-level estimate remains valid for the resulting number-field section integral at $\mathfrak p$.
Exact equality of arbitrary $p$-adic jet coefficients over the number field is not asserted.
\end{proof}

The tensor level has thus been controlled along the sharp sequence: its logarithm is sublinear in the selected degree. The argument does not control the archimedean norm or height of the interpolated section, the cost of imposing the prescribed translated-family conditions, or the source's selected global row-degree sequence. Those quantities still require their own estimates before a shorter theta-lift chain can be concluded or excluded.
